\chapter{Zwartość i jej charakterystyki}
\label{ch:topog-zwartosc}

\noindent\textbf{Cel rozdziału.}
Zwartość jest własnością pozwalającą zastąpić dowolne pokrycie
otwarte pokryciem skończonym. Jej siła ujawnia się w badaniu
odwzorowań ciągłych (w szczególności osiągania kresów
przez funkcje rzeczywiste), rozdzielania zbiorów
i dowolnych iloczynów.
Po definicji i podstawowych przykładach wyprowadzimy równoważne
kryteria przez rodziny zbiorów domkniętych, sieci i ultrafiltry.
Udowodnimy twierdzenie Tychonoffa o zwartości iloczynów oraz
uzgodnione wcześniej twierdzenie o zanurzaniu przestrzeni
Tychonoffa w kostkę $[0,1]^I$. Tam, gdzie wykorzystamy
aksjomat wyboru, zaznaczymy to jawnie.

\section{Definicja zwartości i pierwsze przykłady}
\label{sec:topog-zwartosc-definicja}

\begin{definicja}[Pokrycie otwarte i zwartość]
\label{def:topog-zwartosc}
Niech $(X,\tau)$ będzie przestrzenią topologiczną.
\emph{Pokryciem otwartym} $X$ nazywamy rodzinę
$(U_i)_{i\in I}$ zbiorów z $\tau$ taką, że
$X=\bigcup_{i\in I}U_i$.
\emph{Podpokryciem skończonym} tej rodziny jest
podrodzina $(U_i)_{i\in J}$, gdzie $J\subseteq I$ jest
skończony i $X=\bigcup_{i\in J}U_i$.
Przestrzeń $X$ jest \emph{zwarta}, jeżeli każde jej
pokrycie otwarte ma podpokrycie skończone.
Podzbiór $K\subseteq X$ jest \emph{zwarty w $X$}, gdy
przestrzeń $K$ z topologią podprzestrzeni jest zwarta.
Nie zakładamy żadnego aksjomatu oddzielania.
\end{definicja}

\begin{stwierdzenie}[Zwartość podzbioru w języku przestrzeni otaczającej]
\label{prop:topog-zwartosc-wzgledna}
Niech $K\subseteq X$. Wówczas $K$ jest zwarty dokładnie
wtedy, gdy dla każdej rodziny $(U_i)_{i\in I}$ zbiorów
otwartych w $X$ takiej, że $K\subseteq\bigcup_{i\in I}U_i$,
istnieje skończony $J\subseteq I$ spełniający
$K\subseteq\bigcup_{i\in J}U_i$.
\end{stwierdzenie}
\begin{proof}
Jeśli $K$ jest zwarty, rodzina $(K\cap U_i)$ pokrywa
$K$ zbiorami otwartymi w podprzestrzeni, więc ma
podpokrycie skończone. Odwrotnie, każde pokrycie $K$
zbiorami otwartymi w $K$ ma postać $(K\cap U_i)$
z pewnymi $U_i$ otwartymi w $X$, zgodnie z definicją
topologii podprzestrzeni (rozdział~04). Zastosowanie
warunku dla $U_i$ daje skończone podpokrycie w $K$.
\end{proof}

\begin{przyklad}[Zbiory skończone i przestrzeń pusta]
\label{ex:topog-zwartosc-skonczone}
Każda przestrzeń topologiczna o skończonym zbiorze
punktów jest zwarta. Z dowolnego pokrycia wybieramy
po jednym zbiorze zawierającym każdy punkt; otrzymana
rodzina ma co najwyżej tyle elementów, ile punktów
przestrzeni. Dla $X=\varnothing$ pokrycie puste
ma puste podpokrycie skończone, które pokrywa $X$.
Także każdy skończony podzbiór dowolnej przestrzeni
jest zwarty, niezależnie od jej topologii.
\end{przyklad}

\begin{przyklad}[Niezwartość prostej i przedziału otwartego]
\label{ex:topog-zwartosc-przeciwwklady}
Prosta $\R$ z topologią euklidesową nie jest zwarta:
rodzina $\{(-n,n):n\ge1\}$ pokrywa $\R$, lecz żadna
jej podrodzina skończona tego nie robi, ponieważ
zawiera się w $(-N,N)$ dla największego wybranego $N$.
Podobnie pokrycie $\{(1/n,1):n\ge2\}$ przedziału
$(0,1)$ nie ma podpokrycia skończonego: skończona
rodzina nie pokrywa punktów dostatecznie bliskich $0$.
\end{przyklad}

\begin{figure}[htbp]
\centering
\begin{tikzpicture}[x=9cm,y=.72cm,>=Stealth]
  \draw[->,thick] (-.04,0)--(1.11,0);
  \draw[gray,dashed] (0,-.2)--(0,4.1);
  \draw[gray,dashed] (1,-.2)--(1,4.1);
  \node[below] at (0,-.05) {$0$};
  \node[below] at (1,-.05) {$1$};
  \foreach \a/\y/\n in {.5/1/2,.3333/1.75/3,.25/2.5/4,.1667/3.25/6} {
    \draw[blue!75!black,line width=1.1pt] (\a,\y)--(1,\y);
    \draw[blue!75!black,fill=white,line width=.8pt] (\a,\y) circle (.035cm);
    \draw[blue!75!black,fill=white,line width=.8pt] (1,\y) circle (.035cm);
    \node[anchor=east] at (.015,\y) {$U_{\n}$};
  }
  \node[align=center,blue!55!black,font=\small] at (.61,4.05)
    {Coraz większe zbiory pokrywają $(0,1)$,\\ ale żaden skończony wybór nie wystarcza};
\end{tikzpicture}
\caption{Pierwsze zbiory z pokrycia $U_n=(1/n,1)$, $n\ge2$.
W każdej skończonej podrodzinie pozostaje niepokryty
przedział przy punkcie $0$. Puste końce oznaczają
brak krańców przedziałów w $U_n$.}
\label{fig:topog-zwartosc-pokrycie-otwarte}
\end{figure}

\begin{twierdzenie}[Zwartość przedziału domkniętego]
\label{thm:topog-zwartosc-odcinek}
Dla każdych $a,b\in\R$ z $a\le b$ przedział $[a,b]$
jest zwarty w topologii euklidesowej.
\end{twierdzenie}
\begin{proof}
Przy $a=b$ stosujemy przykład
\ref{ex:topog-zwartosc-skonczone}. Niech $a<b$
i niech $(U_i)_{i\in I}$ będzie pokryciem
otwartym przedziału $[a,b]$ zbiorami otwartymi
w $\R$ (wystarczy dzięki stwierdzeniu
\ref{prop:topog-zwartosc-wzgledna}).
Niech $S$ oznacza zbiór $t\in[a,b]$, dla których
$[a,t]$ pokrywa skończona liczba zbiorów $U_i$.
Zbiór $S$ jest niepusty, ponieważ istnieje $i$
z $a\in U_i$, a wtedy $a\in S$.
Jest też ograniczony z góry przez $b$; niech $c=\sup S$.

Wybierzmy $i_0$ z $c\in U_{i_0}$. Ponieważ $U_{i_0}$
jest otwarty, istnieje $\delta>0$ takie, że
$(c-\delta,c+\delta)\subseteq U_{i_0}$.
Z definicji supremum istnieje $t\in S$ z
$t>c-\delta/2$ (gdy $c=a$, możemy wziąć $t=a$).
Wtedy $[a,t]$ ma pokrycie skończone, a dodanie
$U_{i_0}$ pokrywa także odcinek
$[t,\min\{b,c+\delta/2\}]$.
Gdyby $c<b$, punkt $\min\{b,c+\delta/2\}$
byłby większy od $c$ i należał do $S$, wbrew
definicji supremum. Stąd $c=b$. Ponownie,
ponieważ $t> b-\delta/2$ i $t\in S$, dołączenie
$U_{i_0}$ do skończonego pokrycia $[a,t]$
pokrywa całe $[a,b]$. Dowód jest zakończony.
\end{proof}

\section{Zbiory domknięte, obrazy ciągłe i oddzielanie}
\label{sec:topog-zwartosc-podstawowe-wlasnosci}

\begin{twierdzenie}[Kryterium skończonych przecięć]
\label{thm:topog-zwartosc-fip}
Przestrzeń topologiczna $X$ jest zwarta wtedy i tylko
wtedy, gdy każda \emph{niepusta} rodzina $(F_i)_{i\in I}$
zbiorów domkniętych w $X$, dla której każde
przecięcie skończenie wielu jej elementów jest
niepuste, spełnia $\bigcap_{i\in I}F_i\ne\varnothing$.
Warunek dotyczący skończonych przecięć nazywamy
\emph{własnością skończonych przecięć}.
\end{twierdzenie}
\begin{proof}
Niech $X$ będzie zwarta i załóżmy, że
$\bigcap_i F_i=\varnothing$. Dopełnienia $X\setminus F_i$
tworzą pokrycie otwarte $X$. Istnieje skończona
podrodzina pokrywająca $X$, a z praw de Morgana
przecięcie odpowiednich $F_i$ jest puste,
sprzecznie z założeniem.

Odwrotnie, załóżmy kryterium i niech $(U_i)_{i\in I}$
będzie pokryciem otwartym bez podpokrycia skończonego.
Jeśli $I$ byłby pusty, to $X=\varnothing$, a pusta
podrodzina stanowiłaby podpokrycie skończone.
Zatem $I\ne\varnothing$. Dla $F_i=X\setminus U_i$
każde skończone przecięcie jest niepuste,
gdyż żadne skończenie wiele $U_i$ nie pokrywa $X$.
Z założenia $\bigcap_i F_i\ne\varnothing$, co przeczy
temu, że $(U_i)$ pokrywa $X$.
\end{proof}

\begin{twierdzenie}[Trwałość zwartości]
\label{thm:topog-zwartosc-trwalosc}
Niech $X,Y$ będą przestrzeniami topologicznymi.
\begin{enumerate}[label=(\roman*)]
\item Jeśli $X$ jest zwarta i $F\subseteq X$ jest
domknięty, to $F$ jest zwarty.
\item Jeśli $K\subseteq X$ jest zwarty i
$f\colon X\to Y$ jest ciągłe, to $f(K)$ jest zwarty.
W szczególności ciągły obraz przestrzeni zwartej
jest zwarty.
\end{enumerate}
\end{twierdzenie}
\begin{proof}
Dla (i) niech $(V_i)$ będzie otwartym pokryciem
$F$ w topologii względnej; zapiszmy
$V_i=F\cap U_i$, gdzie $U_i$ są otwarte w $X$.
Rodzina $(U_i)$ uzupełniona o otwarty zbiór
$X\setminus F$ pokrywa $X$. Skończone podpokrycie
$X$ po ograniczeniu do $F$ daje skończone
podpokrycie $(V_i)$ zbioru $F$.

Dla (ii) niech $(W_i)$ będzie otwartym pokryciem
$f(K)$ w $Y$. Ciągłość odwzorowania $f|_K$
sprawia, że $(K\cap f^{-1}(W_i))$ jest otwartym
pokryciem $K$. Wybór skończonego podpokrycia
przenosi się przez $f$ na skończone pokrycie $f(K)$.
\end{proof}

\begin{twierdzenie}[Zwarte podzbiory przestrzeni Hausdorffa]
\label{thm:topog-zwartosc-hausdorff-domkniete}
Jeśli $X$ jest przestrzenią Hausdorffa, to każdy
zwarty podzbiór $K\subseteq X$ jest domknięty w $X$.
\end{twierdzenie}
\begin{proof}
Dla $K=\varnothing$ teza jest oczywista.
Niech $x\in X\setminus K$. Rozważmy wszystkie pary
$(U,V)$ rozłącznych zbiorów otwartych w $X$,
dla których $x\in U$. Zbiory $V$ występujące
w tych parach pokrywają $K$: dla każdego $y\in K$
aksjomat Hausdorffa zapewnia taką parę z $y\in V$.
Ze zwartości $K$ wybieramy skończenie wiele par
$(U_1,V_1),\ldots,(U_m,V_m)$, dla których
$K\subseteq V_1\cup\cdots\cup V_m$.
Zbiór $U=U_1\cap\cdots\cap U_m$ jest
otwartym otoczeniem $x$ i nie przecina $K$:
każdy punkt $K$ należy do pewnego $V_j$,
rozłącznego z $U\subseteq U_j$.
Oznacza to, że $X\setminus K$ jest otwarty,
a więc $K$ domknięty. Konstrukcja nie wymaga
jednoczesnego wyboru otoczeń dla wszystkich $y\in K$.
\end{proof}

\begin{przyklad}[Dlaczego założenie Hausdorffa jest potrzebne]
\label{ex:topog-zwartosc-niedomkniety}
W przestrzeni Sierpińskiego
$X=\{0,1\}$ z topologią
$\{\varnothing,\{1\},X\}$ podzbiór $\{1\}$
jest zwarty, ponieważ jest skończony,
lecz nie jest domknięty: jego dopełnienie $\{0\}$
nie jest otwarte. Nie można więc bez założenia
Hausdorffa wnioskować o domkniętości zbioru zwartego.
\end{przyklad}

\begin{twierdzenie}[Twierdzenie Weierstrassa o osiąganiu kresów]
\label{thm:topog-weierstrass-kresy}
Niech $K$ będzie \emph{niepustą} zwartą przestrzenią
topologiczną, a $f\colon K\to\R$ funkcją ciągłą.
Wówczas $f$ jest ograniczona oraz istnieją punkty
$x_{\min},x_{\max}\in K$ takie, że
\[
 f(x_{\min})=\inf_{x\in K} f(x),\qquad
 f(x_{\max})=\sup_{x\in K} f(x).
\]
W szczególności funkcja $f$ osiąga minimum i maksimum.
Nie zakładamy, że $K$ jest przestrzenią Hausdorffa.
\end{twierdzenie}
\begin{proof}
Na mocy twierdzenia \ref{thm:topog-zwartosc-trwalosc}(ii)
zbiór $A=f(K)$ jest zwarty w $\R$; jest niepusty,
ponieważ $K\ne\varnothing$. Rodzina przedziałów
$((-n,n))_{n\ge1}$ jest otwartym pokryciem $\R$,
a więc również $A$. Ze zwartości $A$ istnieje
$n\ge1$ takie, że $A\subseteq(-n,n)$: jeśli wybrano
skończenie wiele przedziałów, wystarczy wziąć
największy z ich indeksów. Funkcja $f$ jest zatem
ograniczona.

Z zupełności porządku liczb rzeczywistych istnieją
skończone liczby $m=\inf A$ oraz $M=\sup A$.
Ponieważ $\R$ jest przestrzenią Hausdorffa, zbiór
$A$ jest domknięty na mocy twierdzenia
\ref{thm:topog-zwartosc-hausdorff-domkniete}.
Wykażemy, że $M\in A$. Dla każdego $\varepsilon>0$
z definicji supremum istnieje $a\in A$ spełniające
$M-\varepsilon<a\le M$. Stąd każde otwarte otoczenie
$M$ przecina $A$, czyli $M\in\overline A=A$.
Analogicznie dla każdego $\varepsilon>0$ istnieje
$a\in A$ takie, że $m\le a<m+\varepsilon$,
więc $m\in\overline A=A$. Z definicji obrazu
istnieją $x_{\min},x_{\max}\in K$ z
$f(x_{\min})=m$ i $f(x_{\max})=M$.
\end{proof}

\begin{przyklad}[Rola zwartości i ciągłości w twierdzeniu Weierstrassa]
\label{ex:topog-weierstrass-zalozenia}
Funkcja $f\colon(0,1)\to\R$, $f(x)=x$, jest ciągła
i ograniczona, lecz nie osiąga ani swojego infimum $0$,
ani supremum $1$. Jej dziedzina nie jest zwarta.

Z kolei funkcja $g\colon[0,1]\to\R$ określona przez
$g(x)=x$ dla $x<1$ oraz $g(1)=0$ nie jest ciągła
w punkcie $1$. Mimo że $[0,1]$ jest zwarty,
$\sup g([0,1])=1$ nie jest wartością funkcji $g$.
Oba przykłady pokazują, że założeń twierdzenia
nie można po prostu pominąć.
\end{przyklad}

\begin{twierdzenie}[Zwartość i normalność]
\label{thm:topog-zwartosc-hausdorff-normalna}
Każda zwarta przestrzeń Hausdorffa jest $T_4$,
a więc spełnia aksjomat Tychonoffa.
\end{twierdzenie}
\begin{proof}
Przestrzeń Hausdorffa jest $T_1$ (rozdział~05).
Wystarczy dowieść normalności. Niech $F,G\subseteq X$
będą rozłącznymi zbiorami domkniętymi.
Na mocy twierdzenia \ref{thm:topog-zwartosc-trwalosc}
oba są zwarte. Jeśli któryś z nich jest pusty,
możemy użyć $\varnothing$ i $X$ jako rozłącznych
otwartych otoczeń $F$ i $G$. W dalszej części
zakładamy więc, że $F\ne\varnothing$ i $G\ne\varnothing$.

Ustalmy $x\in F$ i rozważmy wszystkie pary
$(U,V)$ rozłącznych zbiorów otwartych takich,
że $x\in U$. Z aksjomatu Hausdorffa ich drugie
składowe $V$ tworzą pokrycie otwarte $G$.
Ze zwartości $G$ wybieramy skończenie wiele par
$(U_1,V_1),\ldots,(U_m,V_m)$, dla których
$G\subseteq V_1\cup\cdots\cup V_m$.
Zbiory $U_x=\bigcap_{j=1}^m U_j$ oraz
$V_x=\bigcup_{j=1}^m V_j$ są otwarte i rozłączne,
$x\in U_x$ oraz $G\subseteq V_x$.

Zatem dla każdego $x\in F$ istnieje para
rozłącznych zbiorów otwartych $(U,V)$ spełniająca
$x\in U$ i $G\subseteq V$. Rozważmy rodzinę
\emph{wszystkich} takich par: ich pierwsze składowe
pokrywają $F$. Zwartość $F$ daje skończenie wiele
par $(U'_1,V'_1),\ldots,(U'_n,V'_n)$ takich,
że $F\subseteq\bigcup_{k=1}^n U'_k$.
Wtedy $U=\bigcup_{k=1}^n U'_k$ i
$V=\bigcap_{k=1}^n V'_k$ są rozłącznymi
otwartymi zbiorami zawierającymi odpowiednio $F,G$.
Otrzymaliśmy normalność; połączenie z $T_1$ daje $T_4$.
Dzięki rozważaniu wszystkich dopuszczalnych par
nie wybieraliśmy jednocześnie otoczeń z nieskończonych
rodzin.
Implikację $T_4\Rightarrow T_{3\frac12}$
udowodniono w rozdziale~05.
\end{proof}

\begin{wniosek}[Ciągła bijekcja z przestrzeni zwartej]
\label{cor:topog-zwartosc-bijekcja-homeomorfizm}
Niech $X$ będzie przestrzenią zwartą, $Y$ przestrzenią
Hausdorffa, a $f\colon X\to Y$ ciągłą bijekcją.
Wówczas $f$ jest homeomorfizmem.
\end{wniosek}
\begin{proof}
Każdy domknięty $F\subseteq X$ jest zwarty na mocy
twierdzenia \ref{thm:topog-zwartosc-trwalosc}.
Jego obraz $f(F)$ jest zwarty jako obraz ciągły,
a więc domknięty w $Y$ na mocy twierdzenia
\ref{thm:topog-zwartosc-hausdorff-domkniete}.
Zatem $f$ jest odwzorowaniem domkniętym.
Bijektywność sprawia, że dla dowolnego domkniętego
$F\subseteq X$ zachodzi $(f^{-1})^{-1}(F)=f(F)$;
stąd odwrotność jest ciągła, więc $f$ jest
homeomorfizmem.
\end{proof}

\section{Zwartość przez sieci i filtry}
\label{sec:topog-zwartosc-zbieznosc}

\begin{definicja}[Punkt skupienia sieci i filtru]
\label{def:topog-zwartosc-punkt-skupienia}
Niech $(x_d)_{d\in D}$ będzie siecią w $X$.
Punkt $x\in X$ nazywamy \emph{punktem skupienia sieci},
jeżeli dla każdego otwartego otoczenia $U$ punktu $x$
i każdego $d_0\in D$ istnieje $d\succeq d_0$
z $x_d\in U$.
Jeśli $\mathcal F$ jest filtrem na $X$, punkt $x$
jest \emph{punktem skupienia filtru}, gdy
$F\cap U\ne\varnothing$ dla każdego
$F\in\mathcal F$ i każdego otwartego otoczenia $U$
punktu $x$. Jest to równoważne
$x\in\bigcap_{F\in\mathcal F}\overline F$,
ponieważ $x\in\overline F$ dokładnie wtedy,
gdy każde otoczenie $x$ przecina $F$.
\end{definicja}

\begin{stwierdzenie}[Punkt skupienia i podsieć zbieżna]
\label{prop:topog-zwartosc-podsiec}
Punkt $x$ jest punktem skupienia sieci $(x_d)_{d\in D}$
wtedy i tylko wtedy, gdy istnieje podsieć $(x_{\varphi(e)})_{e\in E}$
zbieżna do $x$. Przyjmujemy definicję podsieci
z rozdziału~06: odwzorowanie indeksów
$\varphi\colon E\to D$ spełnia warunek ostatecznej
współkońcowości (nie wymagamy jego monotoniczności).
\end{stwierdzenie}
\begin{proof}
Jeżeli podsieć jest zbieżna do $x$, to dla danego
otwartego otoczenia $U$ punktu $x$ od pewnego
$e_0$ mamy $x_{\varphi(e)}\in U$ dla wszystkich
$e\succeq e_0$. Dla dowolnego $d_0\in D$
na mocy kofinalności ogonowej istnieje $e_1$,
po którym $\varphi(e)\succeq d_0$.
Wspólne ograniczenie górne $e_0,e_1$ daje
$e$ z $\varphi(e)\succeq d_0$ i
$x_{\varphi(e)}\in U$, więc $x$ jest punktem skupienia.

Odwrotnie, niech $x$ będzie punktem skupienia.
Oznaczmy przez $\mathcal N_x^o$ rodzinę otwartych
otoczeń $x$ i zbudujmy zbiór
\[
E=\{(d,U,t):d,t\in D,\ d\preceq t,\
U\in\mathcal N_x^o,\ x_t\in U\}.
\]
Z warunku skupienia wynika, że $E\ne\varnothing$.
Definiujemy
$(d,U,t)\preceq_E(d',U',t')$ wtedy i tylko wtedy,
gdy $d\preceq d'$, $U'\subseteq U$ oraz $t\preceq t'$.
Relacja jest zwrotna i przechodnia. Dla dwóch
indeksów wybierzmy $d'$ większy od ich pierwszych
i trzecich współrzędnych; ze skupienia wybierzmy
$t'\succeq d'$ z $x_{t'}\in U_1\cap U_2$.
Indeks $(d',U_1\cap U_2,t')$ jest wtedy wspólnym
ograniczeniem górnym, więc $E$ jest skierowany.

Niech $\varphi(d,U,t)=t$. Monotoniczność wynika
z definicji $\preceq_E$. Dla dowolnego $d_0\in D$
wybierzmy $t_0\succeq d_0$ z
$x_{t_0}\in X$ (możemy wziąć $t_0=d_0$);
indeks $(d_0,X,t_0)\in E$ ma tę własność,
że powyżej niego $\varphi(d,U,t)=t\succeq t_0\succeq d_0$.
Zatem $\varphi$ jest kofinalne w sensie ogonów.
Wreszcie dla otwartego otoczenia $V$ punktu $x$
wybierzmy indeks $(d,V,t)\in E$ ze skupienia.
Wszystkie późniejsze indeksy $(d',U',t')$
spełniają $U'\subseteq V$, zatem
$x_{\varphi(d',U',t')}=x_{t'}\in U'\subseteq V$.
Podsieć zbiega do $x$.
\end{proof}

\begin{twierdzenie}[Kryterium zwartości przez sieci]
\label{thm:topog-zwartosc-sieci}
Przestrzeń topologiczna $X$ jest zwarta wtedy i tylko
wtedy, gdy każda sieć w $X$ ma punkt skupienia.
Równoważnie: każda sieć ma podsieć zbieżną.
\end{twierdzenie}
\begin{proof}
Niech $X$ będzie zwarta i $(x_d)_{d\in D}$
niech będzie siecią w $X$. Przypuśćmy, że nie ma
punktu skupienia. Wówczas rodzina
\[
\mathcal V=\{U\subseteq X:\ U\text{ otwarty oraz }\
\exists d_U\in D\ \forall d\succeq d_U:\ x_d\notin U\}
\]
pokrywa $X$: dla każdego punktu z negacji warunku
skupienia otrzymujemy pewien zbiór należący do
$\mathcal V$, który ten punkt zawiera.
Zwartość daje skończenie wiele
$U_1,\ldots,U_n\in\mathcal V$ pokrywających $X$.
Dla każdego $j$ istnieje $d_j$, po którym sieć
omija $U_j$. Z kierowalności $D$ wybieramy
$d_0\succeq d_1,\ldots,d_n$.
Wtedy $x_{d_0}$ nie należy do żadnego $U_j$,
wbrew temu, że zbiory te pokrywają $X$.
Dla $X=\varnothing$ nie ma żadnych sieci;
implikacja jest wtedy spełniona automatycznie.

Odwrotnie, niech $(U_i)_{i\in I}$ będzie pokryciem
otwartym bez podpokrycia skończonego.
Zamiast wybierać jednocześnie punkty z wielu
niepustych zbiorów, zbudujemy sieć bez użycia
aksjomatu wyboru. Niech
\[
D=\{(J,x):J\subseteq I\text{ skończony},\
 x\in X\setminus\bigcup_{i\in J}U_i\}.
\]
Zbiór $D$ jest niepusty, ponieważ para
$(\varnothing,x)$ istnieje dla pewnego $x\in X$;
gdyby $X=\varnothing$, pokrycie puste byłoby
podpokryciem skończonym.
Ustalmy $(J,x)\preceq(K,y)$ dokładnie wtedy,
gdy $J\subseteq K$. Jest to relacja zwrotna
i przechodnia. Dwie pary mają wspólne ograniczenie
górne: skończona suma $J\cup K$ nie pokrywa
$X$, więc istnieje $z\in X\setminus
\bigcup_{i\in J\cup K}U_i$ i można wziąć
$(J\cup K,z)$. Tak więc $D$ jest skierowany.
Zdefiniujmy sieć $x_{(J,x)}=x$.
Dla dowolnego $z\in X$ wybierzmy $i\in I$
z $z\in U_i$. Ponieważ żadne skończone
zbiory $U_i$ nie pokrywają $X$, istnieje
$(\{i\},x)\in D$. Powyżej tego indeksu
wszystkie wyrazy sieci leżą poza $U_i$;
zatem $z$ nie jest punktem skupienia.
Otrzymujemy sprzeczność i wnioskujemy zwartość.
Równoważność z podsieciami wynika ze stwierdzenia
\ref{prop:topog-zwartosc-podsiec}.
\end{proof}

\begin{definicja}[Ultrafiltr]
\label{def:topog-zwartosc-ultrafiltr}
\emph{Ultrafiltrem} na $X$ nazywamy filtr właściwy
$\mathcal U$, maksymalny względem zawierania wśród
wszystkich filtrów właściwych na $X$.
Filtr właściwy rozumiemy zgodnie z definicją filtru
z rozdziału~06 (w szczególności nie zawiera zbioru pustego). Na zbiorze pustym filtr właściwy
nie istnieje.
\end{definicja}

\begin{lemat}[Kryterium ultrafiltru]
\label{lem:topog-zwartosc-ultrafiltr-alternatywa}
Filtr właściwy $\mathcal U$ na zbiorze $X$ jest
ultrafiltrem wtedy i tylko wtedy, gdy dla każdego
$B\subseteq X$ dokładnie jeden ze zbiorów
$B$ i $X\setminus B$ należy do $\mathcal U$.
\end{lemat}
\begin{proof}
Oba zbiory nie mogą należeć do właściwego filtru,
bo ich przecięcie jest puste. Jeśli $\mathcal U$
jest maksymalny i $B\notin\mathcal U$, to
musi istnieć $F\in\mathcal U$ z $F\cap B=\varnothing$.
W przeciwnym razie rodzina $\mathcal U\cup\{B\}$
miałaby własność skończonych przecięć i generowałaby
filtr właściwy ściśle większy od $\mathcal U$,
co przeczyłoby maksymalności.
Wtedy $F\subseteq X\setminus B$, zatem
$X\setminus B\in\mathcal U$ z domknięcia filtru
na nadzbiory.

Odwrotnie, jeśli filtr wybiera dla każdego $B$
jedną ze stron podziału $X=B\cup(X\setminus B)$,
to żaden większy właściwy filtr nie może zawierać
$B\notin\mathcal U$, bo zawierałby też
$X\setminus B\in\mathcal U$ i stąd zbiór pusty.
To dowodzi maksymalności.
\end{proof}

\begin{lemat}[Przedłużenie filtru do ultrafiltru]
\label{lem:topog-zwartosc-ultrafiltr-zorn}
Każdy filtr właściwy na zbiorze $X$ jest zawarty
w pewnym ultrafiltrze na $X$. W dowodzie korzystamy
z lematu Zorna, a zatem z zasady wyboru.
\end{lemat}
\begin{proof}
Rozważmy uporządkowaną przez zawieranie rodzinę
wszystkich filtrów właściwych zawierających ustalony
filtr $\mathcal F$. Jest niepusta, gdyż zawiera
$\mathcal F$. Jeśli $\mathcal C$ jest łańcuchem,
to suma rodzin $\mathcal G=\bigcup_{\mathcal H\in\mathcal C}
\mathcal H$ jest filtrem właściwym: dla dwóch zbiorów
należących do $\mathcal G$ ich filtry z łańcucha
są porównywalne, zatem większy z nich zawiera oba
i ich przecięcie; domknięcie na nadzbiory
jest równie bezpośrednie, a $\varnothing$ nie należy
do żadnego z filtrów. Dla pustego łańcucha
górnym ograniczeniem jest $\mathcal F$.
Lemat Zorna daje element maksymalny, czyli ultrafiltr.
\end{proof}

\begin{twierdzenie}[Ultrafiltry charakteryzują zwartość]
\label{thm:topog-zwartosc-ultrafiltry}
Przestrzeń topologiczna $X$ jest zwarta wtedy i tylko
wtedy, gdy każdy ultrafiltr na $X$ jest zbieżny
do przynajmniej jednego punktu $X$.
\end{twierdzenie}
\begin{proof}
Niech $X$ będzie zwarta i $\mathcal U$
ultrafiltrem na $X$. Rodzina
$(\overline F)_{F\in\mathcal U}$ jest niepusta
i ma własność skończonych przecięć:
$F_1\cap\cdots\cap F_n\in\mathcal U$ jest
niepusty i zawiera się w każdym $\overline{F_j}$.
Na mocy twierdzenia \ref{thm:topog-zwartosc-fip}
istnieje $x\in\bigcap_{F\in\mathcal U}\overline F$.
Jeżeli otwarte otoczenie $V$ punktu $x$ nie należałoby
do $\mathcal U$, to z lematu
\ref{lem:topog-zwartosc-ultrafiltr-alternatywa}
$X\setminus V\in\mathcal U$.
Mielibyśmy wtedy $x\in\overline{X\setminus V}$,
co jest niemożliwe, ponieważ otwarty zbiór $V$
nie przecina $X\setminus V$.
Zatem każde otwarte otoczenie $x$ należy do
$\mathcal U$, czyli $\mathcal U\to x$.

Odwrotnie, załóżmy, że $X$ nie jest zwarta.
Istnieje pokrycie otwarte $(U_i)_{i\in I}$
bez podpokrycia skończonego. Dopełnienia
$F_i=X\setminus U_i$ mają własność skończonych
przecięć, więc generują filtr właściwy:
jego elementami są nadzbiory niepustych skończonych
przecięć $F_i$. Lemat
\ref{lem:topog-zwartosc-ultrafiltr-zorn}
daje ultrafiltr $\mathcal U$ zawierający każdy $F_i$.
Gdyby $\mathcal U\to x$, wybralibyśmy $i$ z
$x\in U_i$; wtedy również $U_i\in\mathcal U$,
a $F_i=X\setminus U_i\in\mathcal U$, sprzeczność.
Zatem pewien ultrafiltr nie zbiega, co kończy dowód.
\end{proof}

\begin{uwaga}[Punkt skupienia filtru i granica ultrafiltru]
\label{rem:topog-zwartosc-filtr-zbieznosc}
Każdy filtr na przestrzeni zwartej ma punkt skupienia:
rodzina $(\overline F)_{F\in\mathcal F}$ ma własność
skończonych przecięć, bo przecięcie skończonej liczby
elementów filtru jest niepuste. Po zastosowaniu
kryterium \ref{thm:topog-zwartosc-fip} uzyskujemy
$x\in\bigcap_{F\in\mathcal F}\overline F$.
Nie musi to być granica całego filtru; dla ultrafiltru
punkt skupienia jest jednak granicą na mocy dowodu
poprzedniego twierdzenia. W przestrzeni Hausdorffa
każdy filtr ma co najwyżej jedną granicę, jak
udowodniono w rozdziale~06.
\end{uwaga}

\section{Zwartość iloczynów — twierdzenie Tychonoffa}
\label{sec:topog-zwartosc-iloczyny}

\begin{twierdzenie}[Twierdzenie Tychonoffa o iloczynach]
\label{thm:topog-tychonoff-iloczyny}
Niech $I$ będzie dowolnym zbiorem, a
$((X_i,\tau_i))_{i\in I}$ rodziną przestrzeni
zwartych. Wówczas iloczyn
$X=\prod_{i\in I}X_i$ z topologią produktową
jest zwarty. W dowodzie dla dowolnej rodziny
wykorzystujemy aksjomat wyboru.
\end{twierdzenie}
\begin{proof}
Jeśli $I=\varnothing$, iloczyn jest przestrzenią
jednopunktową i jest zwarty. Jeśli któryś czynnik
jest pusty, iloczyn też jest pusty i jest zwarty.
Rozważmy więc rodzinę niepustych przestrzeni.
Aksjomat wyboru zapewnia, że iloczyn $X$ jest
niepusty.

Niech $\mathcal U$ będzie dowolnym ultrafiltrem
na $X$ i $p_i\colon X\to X_i$ rzutowaniem.
Definiujemy filtr obrazowy
\[
\mathcal U_i=\{B\subseteq X_i:
 p_i^{-1}(B)\in\mathcal U\}.
\]
Jest on ultrafiltrem: dla każdego $B\subseteq X_i$
przeciwobrazy $B$ oraz $X_i\setminus B$ są
wzajemnymi dopełnieniami w $X$, więc na mocy lematu
\ref{lem:topog-zwartosc-ultrafiltr-alternatywa}
dokładnie jeden z nich należy do $\mathcal U$.
Pozostałe własności filtru wynikają z własności
przeciwobrazów i tego, że $p_i^{-1}(X_i)=X$.

Zwartość $X_i$ i twierdzenie
\ref{thm:topog-zwartosc-ultrafiltry} dają punkt
$x_i\in X_i$, do którego zbiega $\mathcal U_i$.
Wybierzmy jednocześnie takie punkty dla wszystkich
$i\in I$ (korzystamy tu z aksjomatu wyboru).
Otrzymujemy $x=(x_i)_{i\in I}\in X$.
Dla bazowego otoczenia $x$ w topologii produktowej
\[
W=\bigcap_{j\in J}p_j^{-1}(V_j),\qquad J\subseteq I
\text{ skończony},\quad x_j\in V_j\in\tau_j,
\]
każdy $V_j$ należy do $\mathcal U_j$, a zatem
$p_j^{-1}(V_j)\in\mathcal U$.
Filtr jest domknięty na skończone przecięcia,
więc $W\in\mathcal U$.
Ponieważ zbiory $W$ tworzą bazę otoczeń $x$,
ultrafiltr $\mathcal U$ zbiega do $x$.
Na mocy twierdzenia \ref{thm:topog-zwartosc-ultrafiltry}
iloczyn $X$ jest zwarty.
\end{proof}

\begin{stwierdzenie}[Przypadek skończonego iloczynu]
\label{prop:topog-zwartosc-iloczyn-skonczony}
Dla $n\ge1$ iloczyn przestrzeni zwartych
$X_1\times\cdots\times X_n$ jest zwarty.
Jeżeli ponadto wszystkie te przestrzenie są
Hausdorffa, ich iloczyn jest Hausdorffa.
\end{stwierdzenie}
\begin{proof}
Zwartość wynika z twierdzenia
\ref{thm:topog-tychonoff-iloczyny}.
Jeśli dwa punkty iloczynu są różne, różnią się
w pewnej współrzędnej $i$. Dla tej współrzędnej
wybieramy rozłączne otwarte otoczenia odpowiednich
punktów w $X_i$. Ich przeciwobrazy przez $p_i$
są rozłącznymi otwartymi otoczeniami punktów
w iloczynie, co dowodzi aksjomatu Hausdorffa.
\end{proof}

\begin{twierdzenie}[Twierdzenie Heinego--Borela]
\label{thm:topog-heine-borel}
Dla każdego $n\ge1$ podzbiór $K\subseteq\R^n$
z topologią euklidesową jest zwarty wtedy i tylko
wtedy, gdy jest domknięty i ograniczony.
\end{twierdzenie}
\begin{proof}
Załóżmy, że $K$ jest zwarty.
Przestrzeń $\R^n$ jest Hausdorffa, więc
$K$ jest domknięty na mocy twierdzenia
\ref{thm:topog-zwartosc-hausdorff-domkniete}.
Zbiory $(-m,m)^n$, $m\ge1$, stanowią pokrycie
$\R^n$, więc pewien z nich zawiera $K$:
z każdego skończonego podpokrycia możemy wybrać
największe $m$. To znaczy, że $K$ jest ograniczony.

Odwrotnie, niech $K$ będzie domknięty i ograniczony.
Istnieje $M>0$ z $K\subseteq[-M,M]^n$.
Każdy przedział $[-M,M]$ jest zwarty na mocy
\ref{thm:topog-zwartosc-odcinek}, a jego
$n$-krotny iloczyn jest zwarty na mocy
\ref{prop:topog-zwartosc-iloczyn-skonczony}.
W topologii euklidesowej $\R^n$ topologia
produktowa współrzędnych jest tą samą topologią:
dla metryki maksimum $d_\infty$ i euklidesowej
$d_2$ zachodzi $d_\infty(x,y)\le d_2(x,y)
\le\sqrt n\,d_\infty(x,y)$, a zatem kule
obu metryk są wzajemnie porównywalne i wyznaczają
tę samą rodzinę zbiorów otwartych. Kule metryki
maksimum są właśnie produktami otwartych
przedziałów o wspólnym promieniu.
Zbiór $K$ jest domknięty w $[-M,M]^n$, zatem
jest zwarty na mocy
\ref{thm:topog-zwartosc-trwalosc}(i).
\end{proof}

\begin{figure}[htbp]
\centering
\begin{tikzpicture}[scale=.87,>=Stealth]
\begin{scope}[shift={(0,0)}]
 \fill[blue!10] (0,0) rectangle (3.15,2.5);
 \draw[blue!70!black,thick] (0,0) rectangle (3.15,2.5);
 \foreach \x in {.63,1.26,1.89,2.52} {
   \draw[blue!32] (\x,0)--(\x,2.5);
 }
 \foreach \y in {.5,1,1.5,2} {
   \draw[blue!32] (0,\y)--(3.15,\y);
 }
 \node[fill=white,inner sep=2pt,font=\small] at (1.57,1.25)
   {$[-M,M]\times[-M,M]$};
 \node[font=\small,align=center] at (1.57,-.48)
   {iloczyn dwóch\\zwartych odcinków};
\end{scope}
\begin{scope}[shift={(5,0)}]
 \fill[blue!10] (0,0) rectangle (3.15,2.5);
 \draw[blue!70!black,thick] (0,0) rectangle (3.15,2.5);
 \filldraw[orange!30,draw=orange!85!black,thick]
   (1.52,1.25) ellipse (0.92 and .73);
 \node[font=\small] at (1.52,1.25) {$K$};
 \node[font=\small,blue!70!black] at (2.32,2.23) {$[-M,M]^2$};
 \node[font=\small,align=center] at (1.57,-.48)
   {domknięty podzbiór\\zwartego prostokąta};
\end{scope}
\end{tikzpicture}
\caption{Dwa kroki dowodu twierdzenia Heinego--Borela:
zwartość prostokąta jako iloczynu zwartych odcinków
oraz trwałość zwartości po przejściu do zbioru domkniętego.
Rysunek ma charakter schematyczny.}
\label{fig:topog-heine-borel}
\end{figure}
\FloatBarrier

\section{Zanurzanie przestrzeni Tychonoffa w kostkę}
\label{sec:topog-zwartosc-kostka}

\begin{twierdzenie}[Zanurzanie w kostkę Tychonoffa]
\label{thm:topog-zwartosc-zanurzenie-kostka}
Dla przestrzeni topologicznej $X$ następujące
warunki są równoważne:
\begin{enumerate}[label=(\roman*)]
\item $X$ spełnia aksjomat Tychonoffa, czyli jest
całkowicie regularna i $T_1$;
\item istnieje zbiór indeksów $I$ oraz zanurzenie
topologiczne $e\colon X\to[0,1]^I$.
\end{enumerate}
Można wziąć $I=C(X,[0,1])$, zbiór wszystkich
ciągłych funkcji $X\to[0,1]$, oraz
$e(x)=(f(x))_{f\in I}$.
\end{twierdzenie}
\begin{proof}
Załóżmy (i). Zbiór $I=C(X,[0,1])$ jest zbiorem,
ponieważ jest podzbiorem zbioru wszystkich funkcji
$[0,1]^X$. Definiujemy
$e(x)=(f(x))_{f\in I}$.
Każda współrzędna $p_f\circ e=f$ jest ciągła,
więc $e$ jest ciągłe na mocy własności uniwersalnej
topologii produktowej z rozdziału~04.

Jeżeli $x\ne y$, to $\{y\}$ jest domknięty,
ponieważ $X$ spełnia $T_1$. Całkowita regularność
daje $f\in I$ z $f(x)=0$ i $f(y)=1$.
Zatem $e(x)\ne e(y)$, czyli $e$ jest różnowartościowe.

Pokażemy, że $e^{-1}\colon e(X)\to X$ jest ciągłe.
Niech $U\subseteq X$ będzie otwarty, a $x\in U$.
Zbiór $F=X\setminus U$ jest domknięty i nie zawiera
$x$, więc całkowita regularność daje $f_x\in I$
z $f_x(x)=0$ i $f_x(F)=\{1\}$, o ile $F\ne\varnothing$;
jeśli $F=\varnothing$, wybieramy po prostu funkcję
stałą równą $0$, co wystarcza w poniższym argumencie.
Zbiór
\[
W_x=p_{f_x}^{-1}([0,1/2))\cap e(X)
\]
jest otwarty w $e(X)$, ponieważ $[0,1/2)$ jest otwarty
w $[0,1]$ i rzutowanie $p_{f_x}$ jest ciągłe.
Zawiera $e(x)$ i spełnia $W_x\subseteq e(U)$:
dla $y\in F$ mamy
$f_x(y)=1$, więc $e(y)\notin W_x$.
Ponieważ $x\in U$ było dowolne, każdy punkt
$e(U)$ ma otwarte w $e(X)$ otoczenie zawarte
w $e(U)$. Stąd $e(U)$ jest otwarty w $e(X)$.
To oznacza ciągłość odwrotności i dowodzi,
że $e$ jest zanurzeniem topologicznym.

Odwrotnie, załóżmy (ii). Przestrzeń $[0,1]$ jest
Hausdorffa i całkowicie regularna: dla punktu
$t$ oraz domkniętego $F\subseteq[0,1]$ z $t\notin F$
można, gdy $F\ne\varnothing$, użyć funkcji
$g(s)=|s-t|/(|s-t|+d(s,F))$, gdzie
$d(s,F)=\inf_{z\in F}|s-z|$.
Funkcja $s\mapsto d(s,F)$ jest ciągła:
dla każdego $z\in F$ nierówność trójkąta daje
$|s-z|\le |s-u|+|u-z|$. Biorąc infimum po
$z\in F$ i zamieniając $s$ z $u$, otrzymujemy
$|d(s,F)-d(u,F)|\le |s-u|$.
Mianownik jest dodatni dla $s\ne t$ przez
$|s-t|>0$, a dla $s=t$ przez $d(t,F)>0$,
ponieważ $F$ jest domknięty i $t\notin F$.
Zatem $g$ jest ciągła, $g(t)=0$ i $g|_F=1$.
Dla $F=\varnothing$ wystarczy $g\equiv0$.

Pokażmy bez powoływania się na nieudowodnioną
regułę o iloczynach, że kostka $Q=[0,1]^I$
jest całkowicie regularna. Niech $q\in Q$ i
$F\subseteq Q$ będzie domknięty, $q\notin F$.
Istnieje bazowe otoczenie
$B=\bigcap_{j\in J}p_j^{-1}(V_j)$ punktu $q$
zawarte w $Q\setminus F$, gdzie $J\subseteq I$
jest skończony, a $V_j\subseteq[0,1]$ są otwarte
i zawierają $q_j$.
Dla każdego $j\in J$ wybierzmy ciągłe
$g_j\colon[0,1]\to[0,1]$ z $g_j(q_j)=0$
oraz $g_j=1$ na $[0,1]\setminus V_j$.
Istnienie wynika z całkowitej regularności odcinka;
przy $V_j=[0,1]$ wystarczy $g_j\equiv0$.
Jeśli $J\ne\varnothing$, zdefiniujmy
$G(z)=\max_{j\in J}g_j(z_j)$.
Każda funkcja $g_j\circ p_j$ jest ciągła,
a maksimum dwóch funkcji ciągłych jest ciągłe
na mocy tożsamości
$\max\{a,b\}=(a+b+|a-b|)/2$.
Przez indukcję otrzymujemy ciągłość $G$.
Ponadto $G(q)=0$ oraz $G|_F=1$,
gdyż każdy punkt poza $B$ narusza przynajmniej
jeden z warunków $z_j\in V_j$.
Jeśli $J=\varnothing$, to $B=Q$ oraz
$F=\varnothing$; wówczas weźmy $G\equiv0$.
Ponadto $Q$ jest $T_1$: dla $q\ne r$ istnieje
współrzędna $i$ z $q_i\ne r_i$, a rozłączne
otoczenia w $[0,1]$ dają rozłączne otoczenia
w $Q$. Zatem $Q$ jest Tychonoffa.
Z twierdzenia o dziedziczeniu z rozdziału~05
każda podprzestrzeń $Q$, w szczególności $e(X)$,
jest Tychonoffa. Homeomorfizm
$X\cong e(X)$ przenosi tę własność na $X$.
\end{proof}

\begin{wniosek}[Zwartość i kostkowe uzwarcenie]
\label{cor:topog-zwartosc-uzwarcenie}
Niech $X$ będzie przestrzenią Tychonoffa i niech
$e\colon X\to Q=[0,1]^I$ będzie zanurzeniem
z poprzedniego twierdzenia. Wówczas
$K=\overline{e(X)}^{\,Q}$ jest zwartą przestrzenią
Hausdorffa, a $e(X)$ jest gęsty w $K$.
Jeśli ponadto $X$ jest zwarta, to $e(X)$ jest
domknięty w $Q$ i $K=e(X)$.
\end{wniosek}
\begin{proof}
Na mocy twierdzenia \ref{thm:topog-tychonoff-iloczyny}
kostka $Q$ jest zwarta. Jest Hausdorffa,
jak wykazano w dowodzie twierdzenia
\ref{thm:topog-zwartosc-zanurzenie-kostka}.
Zbiór $K$ jest domknięty w $Q$, więc zwarty
na mocy twierdzenia \ref{thm:topog-zwartosc-trwalosc}.
Jako podprzestrzeń $Q$ jest Hausdorffa.
Gęstość $e(X)$ w $K$ wynika z definicji domknięcia.
Jeśli $X$ jest zwarta, obraz $e(X)$ jest zwarty
jako obraz ciągły, a więc domknięty w
przestrzeni Hausdorffa $Q$. Stąd $K=e(X)$.
\end{proof}

\begin{przyklad}[Prosty model zanurzenia i domknięcia]
\label{ex:topog-zwartosc-uzwarcenie-odcinka}
Funkcja $e\colon\R\to(0,1)$ dana wzorem
$e(t)=\tfrac12+\tfrac{1}{\pi}\arctan t$
jest homeomorfizmem $\R$ na $(0,1)$:
jest ciągła i ściśle rosnąca, jej granice przy
$t\to\pm\infty$ są równe odpowiednio $0$ i $1$,
a odwrotność $t=\tan(\pi(s-1/2))$ jest ciągła
na $(0,1)$. Jako odwzorowanie $\R\to[0,1]$ funkcja $e$ jest
więc zanurzeniem o niedomkniętym obrazie $(0,1)$.
Jego domknięcie $[0,1]$ jest zwarte,
mimo że $\R$ nie jest zwarta.
\end{przyklad}

\section{Podsumowanie i uwagi o założeniach}
\label{sec:topog-zwartosc-podsumowanie}

\begin{center}
\begin{tabular}{@{}p{.33\textwidth}p{.51\textwidth}@{}}
\hline
\textbf{Charakterystyka} & \textbf{Zakres obowiązywania} \\ \hline
Pokrycia otwarte i przecięcia zbiorów domkniętych &
Wszystkie przestrzenie topologiczne.\\
Punkt skupienia każdej sieci; zbieżna podsieć &
Wszystkie przestrzenie topologiczne.\\
Zbieżność każdego ultrafiltru &
Wszystkie przestrzenie topologiczne; dowód przedłużania
filtru używa lematu Zorna.\\
Osiąganie kresów (Weierstrass) &
Każda ciągła funkcja rzeczywista na niepustej
przestrzeni zwartej (bez aksjomatów oddzielania).\\
Zwartość dowolnego iloczynu &
Dla topologii produktowej; w ogólnym przypadku
stosujemy aksjomat wyboru.\\
Domkniętość każdego zwartego podzbioru &
Jeśli przestrzeń otaczająca jest Hausdorffa.\\
Zanurzenie w $[0,1]^I$ &
Dokładnie dla przestrzeni Tychonoffa ($T_{3\frac12}$).\\
\hline
\end{tabular}
\end{center}

\begin{uwaga}[Granice kryteriów ciągowych]
\label{rem:topog-zwartosc-ciagowe}
W ogólnej przestrzeni topologicznej nie zastępujemy
w definicji zwartości sieci dowolnymi ciągami.
Charakterystyka przez sieci i ultrafiltry jest
bezwarunkowa co do aksjomatów oddzielania,
podczas gdy równoważność zwartości ze zwartością
ciągową wymaga dodatkowych hipotez i zostanie
omówiona w odpowiednim późniejszym rozdziale.
\end{uwaga}
